\section{The first two jets on every admissible ray through the origin}
\label{sp:sec:rays}

The uncolored Tornheim function is not jointly holomorphic at the
origin.  Consequently $\mathcal T(0,0,0;1,1)$ is not a value that can
be assigned independently of a restriction convention.  Borwein and
Dilcher established this explicitly for the family
$\mathcal T(t,t,\tau t;1,1)$ and evaluated its first derivative
\cite[Theorems 2.7 and 3.1]{BorweinDilcher}.  The diagonal value $1/3$ is also
recorded by Romik \cite[Theorem 1.3]{sp:Romik}, who credits the simple
diagonal derivative to Borwein and Dilcher.

We give the asymmetric first jet and an exact second-jet formula for
every ray avoiding the two polar directions.  Let
\[
 \mathcal T(A,B,C)=\mathcal T(A,B,C;1,1),\qquad L=\log(2\pi),
 \qquad W_{a,b,c}(t)=\mathcal T(at,bt,ct).
\]
Each restriction is formed away from the polar set and then continued
as a one-variable function.  It is not the evaluation of a nonexistent
joint Taylor series at $(0,0,0)$.

\begin{theorem}[Asymmetric ray values and quadratic jets]
\label{sp:raytheorem}
Let $a,b,c\in\mathbb C$ satisfy $(a+c)(b+c)\ne0$, and define
\begin{equation}\label{sp:rayB}
 \mathcal B(a,b,c)
 =\frac{c(c^2+ab-a^2-b^2)}{(a+c)(b+c)}.
\end{equation}
Then $W_{a,b,c}$ has a removable singularity at $t=0$, and
\begin{align}
 W_{a,b,c}(0)
   &=\frac14+\frac{c}{12(a+c)}+\frac{c}{12(b+c)},
       \label{sp:ray0}\\
 W_{a,b,c}'(0)
   &=\frac{a+b+2c}{4}L+\mathcal B(a,b,c)\zeta'(-1),
       \label{sp:ray1}\\
 W_{a,b,c}''(0)
   &=\frac{2ab+c(a+b)}4L^2
      -\frac{a^2+b^2+2c(a+b)+2c^2}{2}\zeta''(0)\notag\\
   &\hspace{13mm}+(a+b+c)\mathcal B(a,b,c)\zeta''(-1).
       \label{sp:ray2}
\end{align}
\end{theorem}

The equal-inner-order specialization is especially compact:
\begin{align}
 W_{1,1,\tau}(0)&=\frac{5\tau+3}{12(\tau+1)},\notag\\
 W_{1,1,\tau}'(0)&=\frac{\tau+1}{2}L
             +\frac{\tau(\tau-1)}{\tau+1}\zeta'(-1),\notag\\
 W_{1,1,\tau}''(0)&=\frac{\tau+1}{2}L^2
             -(\tau+1)^2\zeta''(0)
             +\frac{\tau(\tau-1)(\tau+2)}{\tau+1}\zeta''(-1).
 \label{sp:equalray}
\end{align}
The first two lines recover the cited classical result; the third is
the quadratic extension proved here.  In particular,
\begin{equation}\label{sp:diagonal2}
 \left.\frac{d^2}{dt^2}\mathcal T(t,t,t)\right|_{t=0}
    =\log^2(2\pi)-4\zeta''(0).
\end{equation}

\subsection{A local meromorphic decomposition}

Put $F(A,B;x)=\operatorname{Li}_A(e^{-x})
                         \operatorname{Li}_B(e^{-x})$.
Jonqui\`ere's convergent expansion for $0<x<2\pi$ gives
\[
 \operatorname{Li}_A(e^{-x})
  =\Gamma(1-A)x^{A-1}
    +\sum_{j\geq0}\zeta(A-j)\frac{(-x)^j}{j!}
\]
for $A$ near zero; all powers of $x$ use its real logarithm.
Subtract the following six terms:
\begin{align}\label{sp:localsix}
 S(A,B;x)={}&\Gamma(1-A)\Gamma(1-B)x^{A+B-2}
   +\Gamma(1-A)\zeta(B)x^{A-1}\notag\\
 &+\Gamma(1-B)\zeta(A)x^{B-1}
   -\Gamma(1-A)\zeta(B-1)x^A\notag\\
 &-\Gamma(1-B)\zeta(A-1)x^B+\zeta(A)\zeta(B).
\end{align}
On a sufficiently small parameter polydisc the remainder and each
fixed parameter derivative are $O(x^{1-\delta}(1+|\log x|^N))$ for
some $\delta<1$.  Consequently the function
\begin{align}\label{sp:Hdef}
 H(A,B,C)={}&\int_0^1 x^{C-1}\{F(A,B;x)-S(A,B;x)\}\,dx
            +\int_1^\infty x^{C-1}F(A,B;x)\,dx\notag\\
 &+\frac{\Gamma(1-A)\Gamma(1-B)}{A+B+C-2}
   +\frac{\Gamma(1-A)\zeta(B)}{A+C-1}
   +\frac{\Gamma(1-B)\zeta(A)}{B+C-1}
\end{align}
is jointly holomorphic near $(0,0,0)$ and symmetric in $A,B$.
The Mellin formula now gives the exact meromorphic decomposition
\begin{align}\label{sp:meromorphicH}
 \mathcal T(A,B,C)=\frac1{\Gamma(1+C)}\biggl\{
  &\zeta(A)\zeta(B)
   -\frac{C\Gamma(1-A)\zeta(B-1)}{A+C}\notag\\
  &-\frac{C\Gamma(1-B)\zeta(A-1)}{B+C}
   +C H(A,B,C)\biggr\}.
\end{align}
To justify it, integrate each term of \eqref{sp:localsix} explicitly
in a common Mellin convergence region, and then continue; the
subtracted integrals in \eqref{sp:Hdef} supply the required local
continuation.  The only polar hyperplanes meeting the origin are now
displayed.  On an admissible ray, their quotients are constants, so
the apparent one-variable singularity is removable.

\subsection{The diagonal input}

The diagonal identity \eqref{sp:diagonal2} has already been proved in
Theorem~\ref{cub:diagonal} by comparing the local cubic residue with
its Hurwitz Fourier representation. That comparison uses only the
coefficient of the simple pole: the unknown cubic finite part first
enters one order later. There is therefore no dependence on a numerical
value of the cubic integral in the argument below.

\subsection{Completing the asymmetric calculation}

The exact restriction
\begin{equation}\label{sp:axis}
 \mathcal T(0,0,C)=\zeta(C-1)-\zeta(C)
\end{equation}
follows by counting the $k-1$ positive pairs summing to $k$ in its
absolute-convergence half-plane.  Restricting
\eqref{sp:meromorphicH} to this axis gives
\[
 \Gamma(1+C)\{\zeta(C-1)-\zeta(C)\}
       =\frac5{12}+C H(0,0,C).
\]
Hence, writing all partial derivatives of $H$ below at the origin,
\begin{align}
 H(0,0,0)
   &=\zeta'(-1)+\frac L2-\frac{5\gamma}{12},\label{sp:H0}\\
 H_C
   &=\frac{\zeta''(-1)-\zeta''(0)}2
       -\gamma\left(\zeta'(-1)+\frac L2\right)
       +\frac5{24}\{\gamma^2+\zeta(2)\}.
       \label{sp:HC}
\end{align}
The symmetry $H_A=H_B$, together with \eqref{sp:diagonal2} in
\eqref{sp:meromorphicH}, then gives
\begin{equation}\label{sp:HA}
 H_A=H_B=\frac{L^2}{8}-\frac{\zeta''(0)}2
         +\gamma\zeta'(-1)-\frac{\gamma L}{4}
         -\frac{\gamma^2+\zeta(2)}{24}.
\end{equation}
These equations determine the entire constant and linear Taylor
polynomial of the jointly holomorphic remainder.

To prove Theorem~\ref{sp:raytheorem}, substitute
$A=at,B=bt,C=ct$ in \eqref{sp:meromorphicH}, use
\eqref{sp:H0}--\eqref{sp:HA}, and expand to second order.  The
constant is \eqref{sp:ray0}.  The coefficient of $\zeta'(-1)$ in
the first derivative is
\[
 c-\frac{bc}{a+c}-\frac{ac}{b+c}
       =\mathcal B(a,b,c),
\]
and all Euler-constant terms cancel, proving \eqref{sp:ray1}.
In the second derivative, the surviving coefficients of
$L^2,\zeta''(0),\zeta''(-1)$ are respectively
\[
 \frac{2ab+c(a+b)}4,\qquad
 -\frac{a^2+b^2+2c(a+b)+2c^2}{2},\qquad
 (a+b+c)\mathcal B(a,b,c).
\]
The coefficients of $\gamma^2$, $\zeta(2)$,
$\gamma\zeta'(-1)$ and $\gamma L$ cancel identically.  This proves
\eqref{sp:ray2} and completes the theorem.

For example, the diagonal and the third-order-parameter axis give
\[
 W_{1,1,1}(0)=\frac13,\qquad W_{0,0,1}(0)=\frac5{12}.
\]
For the asymmetric ray $(a,b,c)=(1,2,3)$ one obtains
\begin{align*}
 W_{1,2,3}(0)&=\frac{29}{80},\\
 W_{1,2,3}'(0)&=\frac94L+\frac9{10}\zeta'(-1),\\
 W_{1,2,3}''(0)&=\frac{13}{4}L^2-\frac{41}{2}\zeta''(0)
                             +\frac{27}{5}\zeta''(-1).
\end{align*}
These are explicit directional identities, and also a concrete audit
test for any routine that substitutes integer orders before
differentiating.  A routine returning the same value for the first
two restrictions has erased the specified limiting direction.

\paragraph{Independent numerical check.}
The accompanying \texttt{verify\_rays.py} integrates the convergent
Jonqui\`ere series on $(0,1)$ and an incomplete-Gamma double series on
$(1,\infty)$.  It does not evaluate a cubic finite-part formula.
At 65-digit working precision, seven successive symmetric step sizes
and polynomial extrapolation reproduce the first two derivatives on
the rays $(1,1,1)$, $(1,2,3)$, $(2,3,1)$, and $(1,-2,3)$.
The largest reported absolute discrepancy is below $2.6\times10^{-37}$.
These are diagnostics of the analytic proof, with truncation and
floating-point errors not certified by interval arithmetic.
